\chapter{双群与旋量特征标的重建方法}
\label{app:double-groups}

\chapref{chap:rotation-group}已经解释，半整数自旋态在 $2\pi$ 转动下获得负号，因此普通点群不足以作为旋量态的线性表示群。现在把连续旋转群的表示限制到晶体保留的有限点群，从覆盖映射、特征标与正交关系重建双群表，并计算轨道与自旋耦合后的分支。双群表由此成为前面表示论和角动量理论的共同应用。

\section{双群作为二重覆盖}

双群表比普通点群表多出来的最重要符号是 $\bar E$。它不是另一个空间恒等操作，而是“在旋量空间中绕任意轴转 $2\pi$”对应的中心群元。对普通三维向量，$E$ 与 $\bar E$ 投影到同一个空间操作；对半整数自旋，二者的表示矩阵却相差负号。因此读双群特征标表时，第一件事就是看某一行在 $\bar E$ 上的特征标：若等于在 $E$ 上的特征标，该表示是单值的；若整体反号，则是自旋型双值表示。这个检查可以在还没理解整张表之前先把普通轨道表示与真正的旋量表示区分开。


\begin{definition}[纯转动点群的双群]
对纯转动点群 $P\subset SO(3)$，定义其双群为覆盖映射
\[
p:SU(2)\longrightarrow SO(3)
\]
下的原像
\begin{equation}
P^D=p^{-1}(P).
\label{eq:double-point-preimage}
\end{equation}
其中中心元记为
\[
\bar E=-I_2,
\qquad
\bar E^2=E,
\]
它在三维空间中对应 $2\pi$ 转动，因此投影到普通点群后与单位元 $E$ 相同，但在旋量表示中不能与 $E$ 识别。
\end{definition}

\begin{proposition}[双群的阶与不可约表示的中心符号]
若 $|P|=N$，则 $|P^D|=2N$，且
\[
1\longrightarrow\{E,\bar E\}\longrightarrow P^D
\overset{p}{\longrightarrow}P\longrightarrow1
\]
为正合列。双群不可约表示分成两类：若
\[
D(\bar E)=+I,
\]
则表示下降为普通点群表示，称为单值表示；若
\[
D(\bar E)=-I,
\]
则为真正的旋量表示或双值表示。
\end{proposition}
\begin{proof}
普通点群的每个元素 $R$ 有两个提升 $u$ 与 $\bar E u=-u$，给出阶数与正合列。因为 $\bar E$ 是中心元，Schur 引理给出 $D(\bar E)=\lambda I$；$\bar E^2=E$ 又要求 $\lambda^2=1$。当 $\lambda=1$ 时，两个提升的表示矩阵相同，所以表示唯一下降到商群 $P^D/\{E,\bar E\}\cong P$；$\lambda=-1$ 时无法下降。
\end{proof}
\chapref{chap:rotation-group}中整数 $j$ 与半整数 $j$ 的差别正对应这两个情况。

含反演的点群还要给状态附加宇称。因为空间反演与纯转动交换，在常见中心对称双群中可以把旋量表示再乘上一维 $g/u$ 宇称表示，所以同一个旋量不可约表示会出现 $+$/$-$ 或 $g/u$ 两个版本。

\section{\texorpdfstring{$SU(2)$ 特征标公式}{SU(2) 特征标公式}}

\begin{proposition}[$SU(2)$ 表示限制到双群的特征标]
自旋 $j$ 不可约表示中，绕任意轴转角 $\theta$ 的提升具有特征标
\begin{equation}
\chi_j(\theta)=\frac{\sin\bigl((j+\tfrac12)\theta\bigr)}{\sin(\theta/2)}
\label{eq:su2-character-angle}
\end{equation}
在 $\theta\to0$ 极限，$\chi_j(0)=2j+1$；而 $2\pi$ 转动给出
\begin{equation}
\chi_j(2\pi)=(-1)^{2j}(2j+1).
\label{eq:central-double-character}
\end{equation}
\end{proposition}
\begin{proof}
任意转轴可共轭到 $z$ 轴，特征标不变。由\chapref{chap:rotation-group}的 $D^{(j)}(U_z\theta)=\operatorname{diag}(\ee^{-\ii m\theta})$，取迹得到有限等比和
\[
\sum_{m=-j}^j\ee^{-\ii m\theta}
=\ee^{\ii j\theta}\frac{1-\ee^{-\ii(2j+1)\theta}}{1-\ee^{-\ii\theta}}
=\frac{\sin((j+\tfrac12)\theta)}{\sin(\theta/2)}.
\]
分母为零时按连续极限取值；在 $\theta=2\pi$ 时，每个本征值都是 $(-1)^{2j}$，故迹为 \eqnrefs{eq:central-double-character}。
\end{proof}
计算双群每个共轭类时，除了空间转角，还必须指定它的两个提升中取哪一个。对相差 $\bar E$ 的提升有 $\chi_j(\bar E u)=(-1)^{2j}\chi_j(u)$，所以半整数自旋不能只靠 $SO(3)$ 转角区分类。

\begin{proposition}[双群限制的重数公式]
若双群不可约特征标记为 $\chi_\alpha$，那么 $j$ 表示限制后的重数为
\begin{equation}
m_\alpha^{(j)}
=\frac1{|P^D|}\sum_{g\in P^D}
\overline{\chi_\alpha(g)}\,\chi_j(g).
\label{eq:double-group-multiplicity}
\end{equation}
\end{proposition}
\begin{proof}
有限群表示完全可约，故 $\chi_j=\sum_\beta m_\beta^{(j)}\chi_\beta$。与 $\chi_\alpha$ 取特征标内积，并用 $\langle\chi_\alpha,\chi_\beta\rangle=\delta_{\alpha\beta}$，便得到重数。
\end{proof}
这就是从“角动量 $j$”转成“晶体双群能带标签”的一般算法。

\section{三个典型双群}

\subsection{\texorpdfstring{$D_2^D$ 与四元数群}{D2 双群与四元数群}}

\begin{proposition}[$D_2$ 的双群与四元数群]
$D_2^D\cong Q_8$。它有四个一维单值不可约表示和一个二维旋量不可约表示。
\end{proposition}
\begin{proof}
普通 $D_2$ 除单位元外，由绕 $x,y,z$ 的三个 $\pi$ 转动组成。它们在 $SU(2)$ 中的提升为
\[
U_x=-\ii\sigma_x,
\qquad
U_y=-\ii\sigma_y,
\qquad
U_z=-\ii\sigma_z.
\]
利用 Pauli 代数，
\[
U_x^2=U_y^2=U_z^2=-I_2=\bar E,
\]
且
\[
U_xU_y=U_z,
\qquad
U_yU_x=-U_z=\bar E U_z.
\]
于是八个矩阵 $\{\pm I_2,\pm U_x,\pm U_y,\pm U_z\}$ 满足四元数群的乘法。其交换子子群为 $\{\pm I_2\}$，Abel 化为 $C_2\times C_2$，故有四个一维普通表示。上述二维矩阵表示不可约，因为 $\sigma_x$ 和 $\sigma_y$ 没有共同本征线。维数平方和为
\[
4\times1^2+2^2=8.
\]
因此这五个表示已穷尽全部复不可约表示。
\end{proof}
这正是\chapref{chap:rotation-group}讨论“原来四个一维 $D_2$ 轨道表示与自旋 $1/2$ 结合后都进入二维旋量表示”的最简代数模型。

\subsection{轴向群与半整数角动量标签}

\begin{proposition}[循环双群的半整数标签]
纯转动群 $C_n$ 的双群为 $C_{2n}$。若 $u$ 是主轴转角 $2\pi/n$ 的标准提升，则 $u^n=\bar E$。其旋量型一维表示以 $m\in\ZZ+\tfrac12$ 模 $n$ 标记，满足
\[
\rho_m(u)=\exp\left(-\ii m\frac{2\pi}{n}\right),
\qquad \rho_m(\bar E)=\ee^{-2\pi\ii m}=-1.
\]
\end{proposition}
\begin{proof}
由 $SU(2)$ 提升公式 $u=\exp[-\ii(\pi/n)\sigma_z]$，$u^n=-I_2$、$u^{2n}=I_2$，且没有更低的正整数幂等于单位元。因此 $\langle u\rangle=C_{2n}$。循环群的一维表示由 $\rho(u)=\ee^{-2\pi\ii q/(2n)}$，$q=0,\ldots,2n-1$ 给出；要求 $\rho(u^n)=-1$ 等价于 $q$ 为奇数，写 $m=q/2$ 就得到半整数标签。$m$ 与 $m+n$ 对生成元和中心元均给出相同值，故在双群中也是同一表示。
\end{proof}
实际晶体表常把 $m$ 与 $-m$ 成对组织，但纯循环双群本身的复不可约表示仍全部是一维。若再有横向二重轴，且它交换两种不同的主轴本征值，两分量才组成不可约二维块；若涉及镜面，需先确定不正转动在自旋空间中的提升约定。

\subsection{\texorpdfstring{立方双群 $O^D$ 与 $O_h^D$}{立方双群 OD 与 OhD}}

在立方对称中，最常见的旋量不可约表示按凝聚态习惯记为
\[
\Gamma_6\ (2\text{维}),
\qquad
\Gamma_7\ (2\text{维}),
\qquad
\Gamma_8\ (4\text{维}).
\]
\begin{proposition}[立方双群的低自旋分支]
把连续转动群的 $j$ 表示限制到 $O^D$，最前几项为
\begin{align*}
j=\frac12&:\quad \Gamma_6,\\
j=\frac32&:\quad \Gamma_8,\\
j=\frac52&:\quad \Gamma_7\oplus\Gamma_8,\\
j=\frac72&:\quad \Gamma_6\oplus\Gamma_7\oplus\Gamma_8.
\end{align*}
维数分别为 $2,4,6,8$。
\end{proposition}
\begin{proof}
取 $\Gamma_6=D^{(1/2)}\downarrow O^D$、$\Gamma_8=D^{(3/2)}\downarrow O^D$，以及 $\Gamma_7=A_2\otimes\Gamma_6$，其中 $A_2$ 为普通八面体群的一维表示。双群的八类与相应特征标为
\begin{center}
\small
\begin{tabular}{c|rrrrrrrr}
\toprule
类 & $E$ & $\bar E$ & $C_3$ & $\bar E C_3$ & $C_4$ & $\bar E C_4$ & $C_2$ & $C_2'$\\
类大小 & $1$ & $1$ & $8$ & $8$ & $6$ & $6$ & $6$ & $12$\\
\midrule
$\chi_6$ & $2$ & $-2$ & $1$ & $-1$ & $\sqrt2$ & $-\sqrt2$ & $0$ & $0$\\
$\chi_7$ & $2$ & $-2$ & $1$ & $-1$ & $-\sqrt2$ & $\sqrt2$ & $0$ & $0$\\
$\chi_8$ & $4$ & $-4$ & $-1$ & $1$ & $0$ & $0$ & $0$ & $0$\\
$\chi_{5/2}$ & $6$ & $-6$ & $0$ & $0$ & $-\sqrt2$ & $\sqrt2$ & $0$ & $0$\\
$\chi_{7/2}$ & $8$ & $-8$ & $1$ & $-1$ & $0$ & $0$ & $0$ & $0$\\
\bottomrule
\end{tabular}
\end{center}
这里 $C_3,C_4$ 指转角分别为 $2\pi/3,\pi/2$ 的提升；两个 $\pi$ 转动类各包含两种提升。其余数值由 \eqnrefs{eq:su2-character-angle} 与中心符号得到。按类大小加权，可直接算出
\[
\langle\chi_6,\chi_6\rangle
=\langle\chi_7,\chi_7\rangle
=\langle\chi_8,\chi_8\rangle=1,
\qquad \langle\chi_\alpha,\chi_\beta\rangle=0\quad(\alpha\ne\beta),
\]
所以前三行为三个不同不可约表示。普通 $O$ 的五个表示贡献维数平方和 $24$，三个旋量表示贡献 $2^2+2^2+4^2=24$，总和为 $|O^D|=48$，故旋量列表完备。表中逐列有 $\chi_{5/2}=\chi_7+\chi_8$、$\chi_{7/2}=\chi_6+\chi_7+\chi_8$，由特征标确定表示即得到全部分支。
\end{proof}
加入反演以后，每个表示再带 $g/u$ 或 $+/-$ 宇称，例如 $\Gamma_8^+$、$\Gamma_7^-$。

这些分支不是经验表。以 $j=5/2$ 为例，先由 \eqnrefs{eq:su2-character-angle} 在 $O^D$ 各类上计算 $\chi_{5/2}$，再用 \eqnrefs{eq:double-group-multiplicity} 与 $\Gamma_6,\Gamma_7,\Gamma_8$ 逐行取内积，就得到
\[
D^{(5/2)}\downarrow O^D=\Gamma_7\oplus\Gamma_8.
\]
实际能带程序输出的双群标签，本质上就是这一限制与约化过程。

\section{轨道表示与自旋表示的组合}

若不考虑自旋--轨道耦合，一个轨道态属于普通点群不可约表示 $\Gamma_{\rm orb}$，自旋 $1/2$ 属于二维旋量表示 $\Gamma_{1/2}$。考虑自旋--轨道耦合以后，应在双群中分解
\begin{equation}
\Gamma_{\rm orb}\otimes\Gamma_{1/2}
=\bigoplus_\alpha m_\alpha\Gamma_\alpha^D.
\label{eq:orbital-spin-double-product}
\end{equation}
重数仍由特征标乘法与正交公式得到。

\begin{proposition}[立方 $d$ 轨道与自旋的张量积分解]
立方 $O_h$ 中，$d$ 轨道先在普通点群下分成
\[
E_g\oplus T_{2g}.
\]
各自再与自旋 $1/2$ 相乘，在常用双群记号下有
\[
E_g\otimes\Gamma_6^+\cong\Gamma_8^+,
\]
\[
T_{2g}\otimes\Gamma_6^+\cong\Gamma_7^+\oplus\Gamma_8^+.
\]
\end{proposition}
\begin{proof}
张量积的特征标等于两因子特征标的乘积。普通 $O$ 的 $E,T_2$ 在 $E,C_3,C_2,C_4,C_2'$ 五类上的特征标分别为
\[
\chi_E=(2,-1,2,0,0),\qquad \chi_{T_2}=(3,0,-1,-1,1).
\]
与上表中的 $\chi_6$ 相乘后，在双群八类上分别得到
\begin{align*}
\chi_E\chi_6&=(4,-4,-1,1,0,0,0,0)=\chi_8,\\
\chi_{T_2}\chi_6&=(6,-6,0,0,-\sqrt2,\sqrt2,0,0)=\chi_7+\chi_8.
\end{align*}
两个因子均取偶宇称时，张量积仍为偶宇称，由此得到所列带 $+$ 的分支。
\end{proof}
这说明“考虑自旋以后能带标签改变”并不是在原来每条能带旁边简单加一个 $\uparrow/\downarrow$；真正的本征态属于双群不可约表示，轨道和自旋通常已经不能分离。

\section{时间反演与双群的关系}

双群解决空间转动在旋量上的线性表示问题；时间反演是反酉对称，需另外说明其平方以及它对动量的作用。
\begin{theorem}[Kramers 简并与 $PT$ 简并]
设 $A$ 是保持某个波矢空间的反酉算符，且 $AH=HA$、$A^2=-I$。则每个本征态 $\ket\psi$ 都有同能量、线性无关的伙伴 $A\ket\psi$。时间反演 $\Theta$ 在 $\vb k=-\vb k+\vb G$ 的动量点满足这些条件时，得到 Kramers 二重简并。若反演 $P$ 也为对称，$P^2=I$、$P\Theta=\Theta P$，则 $(P\Theta)^2=-I$，并在每个 $\vb k$ 保持该波矢空间，因此简并扩展到一般动量。
\end{theorem}
\begin{proof}
若 $H\ket\psi=E\ket\psi$，则 $HA\ket\psi=AH\ket\psi=EA\ket\psi$，因为厄米算符的本征值 $E$ 为实数。反酉性给出 $\langle A\psi,A\phi\rangle=\langle\phi,\psi\rangle$；取 $\phi=A\psi$，得到
\[
\langle A\psi,A^2\psi\rangle=\langle A\psi,\psi\rangle,
\qquad -\langle A\psi,\psi\rangle=\langle A\psi,\psi\rangle,
\]
故两态正交。$P$ 与 $\Theta$ 都把 $\vb k$ 送到 $-\vb k$，其乘积保持 $\vb k$；交换性和 $P^2=I$ 给出 $(P\Theta)^2=P^2\Theta^2=-I$。
\end{proof}
此简并的来源是反酉对称。实际读能带时，要分别检查双群不可约表示的维数与反酉算符平方对简并的要求。
